%%%%%
%
%% Dans ce fichier, les �ventuels accents sont cod�s
%% en ISO-8859-15 (aka latin9)
%% La ligne suivante donne trois exemples 
%% e accent aigu : �, e accent grave : �, e-dans-l'o : �
\section{Extraction du point de grille le plus proche d'un point X(lat,lon,alt)}

\itemstyle{Op\'erateur~:} \operfmt{Extract\_Closest\_PointLatLonAlt}
\index{Extract\_Closest\_PointLatLonAlt}

%%%%%OBJET%%%%
\itemstyle{Objet~:}
Extraction du point de grille le plus proche d'un point X d�crit par ses coordonn�es latitutde,longitude,altitude.

%%%%%ARGUMENTS%%%%
\bigskip
%% Dans ce fichier, les �ventuels accents sont cod�s
%% en ISO-8859-15 (aka latin9)
%% La ligne suivante donne trois exemples 
%% e accent aigu : �, e accent grave : �, e-dans-l'o : �
%%%%%ENTREES%%%%
\begin{tabarg}
\multicolumn{3}{l}{\titretab {Entr\'ees}}\\
\hline
$X\_in$ & Vecteur, \texttt{float} &  Vecteur d'entr�e\\
$P\_in$ & Vecteur, \texttt{float} & Altitude en pression\\
$tgt\_lat$ & Coeff., \texttt{float} & Latitude du point cible en degr� d�cimal\\
$tgt\_lon$ & Coeff., \texttt{float} & Longitude du point cible en degr� d�cimal\\
$tgt\_alt$ & Coeff., \texttt{float} & Altitude du point cible en degr� d�cimal\\
\end{tabarg}

\medskip
%%%%%SORTIES%%%%
\begin{tabarg}
\multicolumn{3}{l}{\titretab {Sorties}}\\
\hline
$X\_out$ & Vecteur, \texttt{float} & Colonne extraite du vecteur d'entr�e la plus proche du point cible\\
\end{tabarg}


%%%%%FORMULE%%%%
\itemstyle{Formule}

%% Formule de l'operateur : PressAlti
%% Dans ce fichier, les �ventuels accents sont cod�s
%% en ISO-8859-15 (aka latin9)
%% La ligne suivante donne trois exemples 
%% e accent aigu : �, e accent grave : �, e-dans-l'o : �
\[
X\_out = (lats-tgt\_lat)*RAYON\_TERRE)**2 + (lons-tgt\_lon)*RAYON_TERRE*cos(lats)**2+\frac{1}{9.80665*alts-tgt\_alt}**2 
\]
Avec : RAYON\_TERRE = 6375e3



%%%%%AUTEUR%%%%
\itemstyle{Auteur~:}
CNRM/GMEI/TRAMM

%%%%%COMMENTAIRE%%%%
\itemstyle{Commentaire~:}
Les degr�s sont donn�s en degr�s d�cimaux en entr�e de l'op�rateur, qui se charge de convertir en radians.\\

